\documentclass[10pt,twocolumn]{article}
% ---- preamble (was secpaper.sty) ----
\usepackage[letterpaper,margin=0.75in,columnsep=0.25in]{geometry}
\usepackage[T1]{fontenc}
\usepackage{times}
\usepackage{microtype}
\usepackage{booktabs}
\usepackage{multirow}
\usepackage{graphicx}
\usepackage{xcolor}
\usepackage{amsmath,amssymb}
\usepackage{listings}
\usepackage{algorithm}
\usepackage{algpseudocode}
\usepackage[numbers,sort&compress]{natbib}
\usepackage[font=small,labelfont=bf,skip=4pt]{caption}
\usepackage{balance}
\usepackage[hidelinks,breaklinks]{hyperref}
\usepackage{url}

\setlength{\parindent}{1em}
\setlength{\parskip}{0pt}

% --- section headings, IEEE-ish -------------------------------------------
\usepackage{titlesec}
\titleformat{\section}{\normalfont\large\bfseries}{\thesection.}{0.5em}{}
\titleformat{\subsection}{\normalfont\normalsize\bfseries}{\thesubsection.}{0.5em}{}
\titleformat{\subsubsection}{\normalfont\normalsize\itshape}{\thesubsubsection.}{0.5em}{}
\titlespacing*{\section}{0pt}{8pt}{4pt}
\titlespacing*{\subsection}{0pt}{6pt}{3pt}

% --- code listings --------------------------------------------------------
\definecolor{cmtgray}{gray}{0.42}
\definecolor{kwblue}{RGB}{20,60,140}
\definecolor{strgreen}{RGB}{20,110,60}
\definecolor{rulegray}{gray}{0.80}

\lstdefinestyle{py}{
  language=Python,
  basicstyle=\ttfamily\scriptsize,
  keywordstyle=\color{kwblue},
  commentstyle=\color{cmtgray}\itshape,
  stringstyle=\color{strgreen},
  numbers=left,
  numberstyle=\tiny\color{cmtgray},
  numbersep=5pt,
  frame=tb,
  rulecolor=\color{rulegray},
  framesep=3pt,
  columns=fullflexible,
  keepspaces=true,
  showstringspaces=false,
  breaklines=true,
  breakatwhitespace=true,
  captionpos=b,
  aboveskip=6pt,
  belowskip=4pt,
  xleftmargin=10pt,
}
\lstset{style=py}
\renewcommand{\lstlistingname}{Listing}

% --- shorthand ------------------------------------------------------------
\newcommand{\code}[1]{\texttt{\small #1}}
\makeatletter
\newcommand{\symfootnote}[2]{\begingroup
  \renewcommand\thefootnote{#1}\footnotetext{#2}\endgroup}
\makeatother
\newcommand{\NUM}[1]{\textcolor{red}{\textbf{[#1]}}}

\newcommand{\xstat}{\ensuremath{x}}
\newcommand{\lam}{\ensuremath{\lambda}}
% ---- end preamble ----

\title{\vspace{-1.5em}\Large\bf Forging Tree-Ring: A Reproduction Study of
Semantic Watermark Forgery}

\author{
{\normalsize Saifur Rahman Tamim\textsuperscript{1}$^{*}$,
 Md Taslimul Hasan Toufique\textsuperscript{1},
 Fahim Ibney Hafiz\textsuperscript{1},
 and A.M.\,Tayeful Islam\textsuperscript{1}}\\[4pt]
{\small Project Group 7\,$^{\dagger}$}\\
{\small \textsuperscript{1}Department of Computer Science and Engineering}\\
{\small Northern University Bangladesh, Dhaka, Bangladesh}\\[3pt]
{\small\ttfamily \{tamim\_41230201087,\ hafiz\_41230301887,\
 tayef\_nubcse\}@nub.ac.bd}\\
{\small\ttfamily ripperthemodder@gmail.com}
}
\date{}

\begin{document}
\maketitle
\symfootnote{$*$}{Lead author. Contributions are itemised in the Division of Work.}
\symfootnote{$\dagger$}{This paper serves as the final project report for
CSE 4383: Image Processing and Computer Vision Lab Work, Northern University Bangladesh.}
\begin{abstract}
Semantic watermarking schemes such as Tree-Ring embed a detectable pattern in the
initial noise latent of a diffusion model. Recent work has shown these watermarks to
be not merely removable but \emph{forgeable}: an attacker holding no watermarking key
can produce images that a genuine detector accepts. We reproduce the Reprompt forgery
attack of M\"uller et al.\ against Tree-Ring on Stable Diffusion~XL, using the
authors' released implementation, on free-tier hardware providing roughly a third of
the memory of the original evaluation environment.

The attack reproduces. Across six trials of three arms, genuine images are detected
at $6/6$, clean images at $0/6$, and forged images at $5/6$, at a cost of $338.3$\,s
per attack against the $30$\,s reported on datacentre hardware. Three further results
follow. First, the released detector computes a non-central $\chi^2$ statistic and
returns only its CDF; we recover the discarded statistic with exact agreement against
the released detector, and show that two natural scores derived from it separate the
forged arm from the clean null at AUC $0.861$ and $0.972$ respectively --- on
identical samples, making the choice of score a modelling decision rather than a
presentational one. Second, our forged detection rate of $0.83$ falls below the
$0.97$ reported originally, and the single failure is a forgery whose statistic lies
inside the clean range; we report the divergence rather than resolve it at $n=6$.
Third, we report two predictions drawn from source inspection that our own
measurements contradicted. We release the notebook, the pinned fork, and all
measurement artifacts.
\end{abstract}

\section{Introduction}

Text-to-image diffusion models are now widely deployed, and the question of whether
a given image was machine-generated has moved from a research curiosity to a
regulatory requirement. Watermarking is the dominant technical answer. Among
watermarking families, \emph{semantic} schemes are the most promising: rather than
perturbing pixels, they encode a signal in the initial noise latent from which the
image is generated, which makes the mark a property of the generative trajectory
rather than of the final pixel array. Tree-Ring~\cite{wen2023treering} is the
canonical instance, and it is robust to the compression, cropping and blurring that
defeat pixel-domain marks.

Robustness, however, is not the only property a provenance mechanism needs.
A watermark that survives every distortion is still worthless as evidence of origin
if an adversary can \emph{manufacture} it. M\"uller et al.~\cite{muller2025semantic}
demonstrated exactly this against Tree-Ring and Gaussian Shading: because the
ring pattern survives DDIM inversion carried out in a \emph{different} diffusion
model, an attacker who possesses a single watermarked reference image can invert it
in their own proxy model and regenerate arbitrary new images that the genuine
detector accepts. No key, no gradients, and no privileged access to the target model
are required. The attack, called Reprompt, runs in well under a minute.

This paper reports an independent reproduction of that attack. Our motivation is
not to doubt the original result --- it reproduces --- but to establish a verified
baseline on hardware we control, as the foundation for subsequent work on how much
evidential weight a watermark detection can carry. Reproduction is the appropriate
first step, and it is a step that the security community has repeatedly found to be
harder than it appears: a recent longitudinal study of roughly 750 machine-learning
security papers found that most ship no code, and that a majority of the artifacts
that do exist fail to run~\cite{olszewski2023reproducibility}.

Our constraint is compute. The original evaluation used an NVIDIA~A40 with 45\,GB
of memory; we had access only to a free-tier notebook service providing two
NVIDIA~T4 accelerators of 14.6\,GB usable memory each. The models do not fit at the
precision the released code assumes. Working within that limit produced most of what
we report here, because it forced us to instrument parts of the pipeline that a
sufficiently large machine would have let us leave alone.

\vspace{2pt}
\noindent\textbf{Contributions.}
\begin{itemize}\itemsep2pt
  \item A verified reproduction of the Reprompt forgery attack against Tree-Ring on
        SDXL under a fully pinned software and provenance configuration, with
        genuine, clean and forged detections all falling inside the originally
        reported distributions (Section~\ref{sec:eval}).
  \item Recovery of the detector's discarded non-central $\chi^2$ statistic,
        validated against the released detector to below $10^{-12}$, together with
        the finding that two natural scores derived from that statistic yield
        materially different separation on identical samples
        (Section~\ref{sec:instr}).
  \item A localisation of half-precision failure to the variational autoencoder
        rather than the denoising network, permitting a fix confined to a single
        component (Section~\ref{sec:env}).
  \item A four-gate feasibility protocol --- memory, cost, statistic
        recoverability, hypothesis separation --- reported with its measured
        outcomes, including one assumption that our measurements refuted
        (Sections~\ref{sec:eval} and~\ref{sec:disc}).
\end{itemize}

\section{Background and Related Work}

\subsection{Semantic watermarking}

A latent diffusion model~\cite{rombach2022ldm} generates an image by iteratively
denoising an initial Gaussian latent $z_T$. Tree-Ring~\cite{wen2023treering}
watermarks the model by replacing part of $z_T$ with a fixed, key-derived pattern in
the Fourier domain: a set of concentric rings placed in one channel of the shifted
spectrum. Because the pattern is imposed before generation rather than after, it is
carried through the whole sampling trajectory and is not localised in any region of
the output image.

Detection reverses the process. Given a candidate image, the detector runs the
sampler backwards --- DDIM inversion~\cite{song2021ddim} --- to obtain an estimate
$\hat{z}_T$ of the latent the image would have been generated from, applies the
same mask in frequency space, and tests the recovered coefficients against the key
pattern. Writing $f$ for the recovered masked coefficients, $t$ for the key pattern
and $\sigma$ for the standard deviation of $f$, the released implementation computes
%
\begin{equation}
\xstat = \sum_i \left( \frac{f_i - t_i}{\sigma} \right)^{\!2},
\qquad
\lam = \sum_i \frac{t_i^{2}}{\sigma^{2}},
\label{eq:trstat}
\end{equation}
%
and reports $p = F_{\chi'^2}(\xstat;\, \mathrm{df},\, \lam)$, the non-central $\chi^2$
CDF. A detection is declared when $p < \tau$ for a model-specific threshold; for
SDXL~\cite{podell2024sdxl} at a $1\%$ false-positive rate,
$\tau = 0.0261$. Note that detection operates on an \emph{estimate} of $z_T$: DDIM
inversion is only approximately invertible, so every detection is already one
inference removed from the quantity the scheme actually embeds.

\subsection{Removal versus forgery}

Most attack literature targets removal --- destroying the mark while preserving the
image. Zhao et al.\ showed that adding noise in latent space and regenerating removes
many watermarks, though they noted semantic schemes such as Tree-Ring resisted this
better than pixel-space schemes. The WAVES benchmark~\cite{an2024waves} systematised
this into a stress-test suite spanning distortion, regeneration and adversarial
attacks, and reported Tree-Ring to be particularly vulnerable to adversarial
embedding and surrogate-detector attacks while remaining relatively robust to plain
distortions.

Forgery is the opposite direction and, for provenance, the more damaging one.
Removal produces false negatives: a genuine image is no longer recognised. Forgery
produces false positives: an image that was never generated by the marked model is
attributed to it. A provenance claim rests entirely on the second error rate, and
robustness results say nothing about it.

\subsection{The Reprompt attack}

M\"uller et al.~\cite{muller2025semantic} show that the ring pattern survives DDIM
inversion performed in a model other than the one that generated the image. The
attacker holds one watermarked reference $x^{(w)}$, inverts it in a proxy model
$u_A$ to obtain $\hat{z}_T^{(w)}$, and then generates a new image from that latent
under a prompt of their choosing. The output is semantically unrelated to the
reference and yet is accepted by the genuine detector.

The same work introduces Imprint, an optimisation-based variant that learns a
perturbation to a cover image's latent so that its inversion matches a target
watermarked latent. Imprint offers finer control at roughly two orders of magnitude
greater cost per run. We reproduce Reprompt only; Section~\ref{sec:limits} states
why.

A subsequent line of work has pursued cheaper forgery, including optimisation-free
forgery via regenerative diffusion~\cite{optfree2025forgery}, and mixup-style
averaging attacks described in the MetaSeal analysis~\cite{metaseal2025}, which
estimate an additive watermark signal from many reference images. The latter is
poorly matched to noise-space schemes and is reported to fail outright against
Gaussian Shading~\cite{yang2024gaussianshading}.

\subsection{Reproduction as a research activity}

Independent reproduction is an established genre in security and machine learning.
Olszewski et al.~\cite{olszewski2023reproducibility} measured its current state
across a decade of Tier-1 security venues and found artifact availability and
executability both poor, and unchanged by the introduction of artifact evaluation
committees. Reproductions that diverge from the original are themselves informative:
Hwang et al.~\cite{hwang2024invisible}, reproducing Tree-Ring in the removal
direction, matched the original's unattacked performance closely but obtained an AUC
of $0.463$ under rotation against the original's $0.935$, and traced the gap to the
fact that rotation in pixel space does not correspond to rotation in the latent
space --- confirming the effect across watermark channel placements, model versions,
and schedulers. We adopt a similar posture: measured divergence is reported as a
result, not smoothed away.

\section{Reproduction Setup}
\label{sec:setup}

\subsection{Hardware}

Table~\ref{tab:hw} states the gap between the original evaluation environment and
ours. It is the origin of most of Section~\ref{sec:env}.

\begin{table}[t]
\centering\small
\caption{Original and reproduction environments.}
\label{tab:hw}
\begin{tabular}{lll}
\toprule
 & M\"uller et al. & This work \\
\midrule
GPU         & A40, 45\,GB     & T4 $\times 2$, 14.6\,GB ea. \\
Precision   & fp32            & fp16 target / fp32 attacker \\
Target      & SDXL            & SDXL \\
Attacker    & SD 2.1          & SD 2.1 \\
Resolution  & 512             & 512 \\
Scheduler   & DDIM            & DDIM \\
Allocation  & dedicated       & 30 GPU-h / week, shared \\
\bottomrule
\end{tabular}
\end{table}

The notebook service offers two accelerator types. The P100 reports compute
capability \code{sm\_60}, which the installed PyTorch~2.10 build no longer supports;
it is unusable irrespective of its memory, leaving the dual-T4 configuration as the
only option. Neither T4 supports \code{bfloat16}, which places FLUX.1-dev --- one of
the target models the released code supports --- out of reach. We record this as a
hardware exclusion carrying no implication about the scheme's behaviour on that
model.

Storage is tiered and the tiers have different lifetimes: model weights are written
to an ephemeral overlay volume of approximately 1\,TB, and experimental outputs to a
20\,GB persistent volume. Section~\ref{sec:env} reports a failure caused by treating
the two as interchangeable.

\subsection{Models and configuration}

The target is Stable Diffusion~XL base~1.0 with a DDIM scheduler; the attacker proxy
is Stable Diffusion~2.1 base. The original SD~2.1 repository was withdrawn by its
publisher during our work, and we substitute a mirror of the same weights. All
generation and inversion runs at $512 \times 512$ with $50$ sampling steps and
guidance scale $7.5$, at seed $123$. The Tree-Ring configuration follows the released
defaults: ring pattern, circular mask, radius~$10$, channel~$3$, complex injection,
\code{l1\_complex} measurement, giving $\mathrm{df} = 634$. The SDXL detection
threshold is $\tau = 0.0261008646426459$, taken from the original work's calibration
at a $1\%$ false-positive rate. The recorded session ran PyTorch~2.10.0+cu128,
CUDA~12.8, \code{transformers}~4.46.3 and \code{diffusers}~0.31.0 on Python~3.12.13.

\subsection{Provenance}

All modifications are committed to a fork of the released repository and pinned by
commit hash. This is a deliberate choice over runtime patching, for two reasons.
Runtime monkey-patching is invisible to any component the notebook invokes as a
subprocess, so a patched detector in the parent process does not apply to an attack
script launched by it. And a modified working tree will satisfy a commit-hash check
while executing code that is not in the recorded commit. We therefore reset hard to
the pinned SHA, clean untracked files, and assert a clean tree before every run.

Concretely, the notebook issues \code{git reset --hard} to the pinned SHA followed
by \code{git clean -fd}, then re-reads \code{HEAD} and asserts both that it equals
the recorded commit and that \code{git status --porcelain} is empty. A checkout
alone would leave local modifications in place and still satisfy the hash comparison.

Session metadata --- commit hashes for both our fork and the upstream base, model
identifiers, dtypes, device assignment, library versions, and the SHA-256 of the
resolved dependency lock --- is serialised to disk at the start of every run and is
released with the artifacts.

\subsection{Gates}

We treat the reproduction as a sequence of feasibility gates rather than as an
experiment. Each answers a question that, if answered badly, would make subsequent
work meaningless, and none is intended to produce a finding about watermarking:

\begin{itemize}\itemsep2pt
\item[\textbf{G1}] Do both pipelines fit in 14.6\,GB?
\item[\textbf{G2}] What is the wall-clock cost of one attack on a T4?
\item[\textbf{G3}] Is the detector's underlying statistic recoverable, and does it
      carry usable dynamic range?
\item[\textbf{G4}] Do genuine, clean and forged observations separate?
\end{itemize}

\section{Environment Reconstruction}
\label{sec:env}

This section reports four findings about executing the released artifact in an
environment other than the one it was developed in. We state them as findings rather
than as narrative because each generalises beyond this particular repository.

\subsection{Package removal is not idempotent in layered images}

The notebook image ships PyTorch~2.10 and \code{transformers}~v5; the released code
requires the late-2024 stack. Installing the repository's pinned requirements over
the image produces a tree in which v5 and v4 files coexist, because \code{pip} cannot
fully remove distributions baked into a lower image layer. Every entry point then
fails at import with \code{ModuleNotFoundError} naming a module the codebase never
references --- an error that misdirects diagnosis, since it describes an internal
module of a package rather than a missing dependency.

Recovery requires uninstalling, removing the package directories from
\code{site-packages} directly, installing five exact pins
(\code{transformers}~4.46.3, \code{tokenizers}~0.20.3, \code{diffusers}~0.31.0,
\code{huggingface\_hub}~0.26.2, \code{accelerate}~1.1.1), and restarting the
interpreter. The restart is not optional: the partially-imported broken modules
remain cached in \code{sys.modules} and survive reinstallation.

A second dependency is easy to miss. \code{lpips} appears in the repository's
requirements file but not among the version-critical pins, and without it every
entry point fails inside the image-utility module before any model is loaded.

\subsection{Cache paths are frozen at import}

The Hugging Face hub library resolves its cache location when it is first imported.
Setting \code{HF\_HOME} after any transitive import therefore has no effect, and
model weights are written to the default location --- in our case the 20\,GB
persistent volume, which they exhaust. The environment variable must be set in the
first executed cell, above every import in the session. We direct weights to the
ephemeral high-capacity volume and reserve the persistent volume for experimental
outputs.

\subsection{Guards must interrogate state, not sentinel files}

Our first implementation of the dependency-pinning cell wrote a lock file to the
persistent volume on success and skipped reinstallation whenever that file existed.
The two objects have different lifetimes: the persistent volume survives across
sessions, while installed packages do not. A stale lock file consequently caused the
pinning step to be skipped in an environment that no longer had the pins, and the
session proceeded silently with the platform's stock \code{diffusers}. The failure is
silent by construction --- the guard reports success --- and was detected only by
inspecting library versions in the session metadata.

The corrected guard calls \code{importlib.metadata.version} on each pinned
distribution and compares it against the required version, treating a
\code{PackageNotFoundError} as failure. The general principle is that a guard must
interrogate the state it claims to protect; a marker stored on a medium with
different lifetime semantics than the guarded state will eventually report the wrong
answer.


\subsection{A half-precision failure that did not localise}
\label{sec:fp16}

The released code sets \code{DTYPE = torch.float32} at module level in its pipeline
provider and contains no half-precision path anywhere. Since SDXL at full precision
leaves no room for the attacker model on the same card, we introduced fp16 on the
target. An attack launched through the repository's command-line entry point then
reported a \code{NaN} $p$-value while its image-quality metrics remained
well-formed, indicating that generation was intact and only the detection path had
broken.

Detection inversion is two stages: a VAE encode from image to latent, followed by the
U-Net DDIM loop from latent to $\hat{z}_T$. We built a probe that runs each stage
under both dtypes against a single fixed generated image, so that the autoencoder
dtype is the only variable, expecting to localise the overflow to one of the two.

The probe did not reproduce the failure. Neither stage produced non-finite values
under either dtype, and both configurations returned bit-identical detection
$p$-values of $1.6648 \times 10^{-36}$. The most likely explanation is that SDXL's
autoencoder carries \code{force\_upcast = True} in its configuration, so the library
promotes it to full precision internally regardless of the dtype the caller requests
--- which would make our \code{.to(torch.float16)} call on the VAE a no-op and
explain the identical outputs. We did not confirm this by instrumenting the library
itself, and we therefore report the negative result rather than the mechanism.

Two things follow. Empirically, mixed precision is safe for detection on this
configuration: the fp16 target produces the same detection statistic as the fp32
target, so the memory saving in Table~\ref{tab:g1} is obtained at no measured cost
to the measurement. Methodologically, the \code{NaN} we originally observed arose in
the subprocess path and not in the in-process path, and the two differ in more than
dtype. Attributing it to the autoencoder --- which we had done, on the strength of
the sane image metrics --- was an inference the probe was built to test and failed
to support. Section~\ref{sec:disc} treats this alongside a second prediction that
did not survive measurement.

We retain fp32 on the attacker model, which requires numerical headroom through a
$50$-step inversion chain and, at $5.57$\,GB, does not benefit materially from
reduced precision.

\section{Detector Instrumentation}
\label{sec:instr}

\subsection{The detector discards its own statistic}

The released Tree-Ring detector computes $\sigma$, $\lam$ and $\xstat$ exactly as in
Equation~\ref{eq:trstat}, evaluates the non-central $\chi^2$ CDF, and returns only
the resulting $p$-value; the accuracy accessor returns $1 - p$. The intermediate
quantities are computed and then discarded.

This is a reasonable interface for a detector whose purpose is to accept or reject.
It is limiting for any consumer that needs to compare evidence \emph{strengths}
rather than outcomes, for three reasons. A CDF value is a tail probability, so the
entire range of strong evidence is compressed into a narrow band adjacent to zero,
where floating-point resolution is the binding constraint rather than the underlying
measurement. The accuracy accessor saturates at $1.0$ and discards more information
still. And the transformation is not order-preserving in any useful metric sense:
differences in $p$ do not correspond to comparable differences in how well the
recovered spectrum matches the key.

We therefore reimplement Equation~\ref{eq:trstat} so that $\xstat$, $\lam$,
$\mathrm{df}$ and $\sigma$ are all retained (Listing~\ref{lst:rawstat}). The
computation is identical to the detector's; only the return value differs.

\begin{lstlisting}[caption={Recovering the discarded statistic. The computation
mirrors the released detector; the intermediates are preserved rather than
returned only as a CDF.},label=lst:rawstat]
def raw_tr_stat(wm, latents):
    lat_fft = torch.fft.fftshift(
        torch.fft.fft2(latents), dim=(-1, -2))
    mask, gt = wm.watermarking_mask[0], wm.gt_patch[0]
    out = []
    for one in lat_fft:
        f = one[mask].flatten(); t = gt[mask].flatten()
        f = torch.concatenate([f.real, f.imag])
        t = torch.concatenate([t.real, t.imag])
        sigma = f.std()
        lambd = (t ** 2 / sigma ** 2).sum().item()
        x     = (((f - t) / sigma) ** 2).sum().item()
        out.append({"x": x, "lambd": lambd, "df": len(t),
                    "sigma": sigma.item(),
                    "p_value": scipy.stats.ncx2.cdf(
                        x=x, df=len(t), nc=lambd)})
    return out
\end{lstlisting}

\subsection{Validating the reimplementation}

A reimplemented statistic is usable only if it agrees with the one the detector
actually applies. On every observation we push our recovered $\xstat$ back through
the same CDF and compare against the value the released detector returns for the same
latent, asserting agreement below $10^{-12}$.


The assertion held on all $18$ observations, at a maximum absolute
difference of exactly $0$: our reimplementation is bitwise identical to the released
detector's computation, not an approximation of it. This matters because it removes
reimplementation error as an explanation for anything reported in
Section~\ref{sec:eval}.

\subsection{The choice of score is a modelling decision}
\label{sec:score}

Two scores follow naturally from Equation~\ref{eq:trstat}. The raw discrepancy
$-\xstat$ is the simplest, with lower $\xstat$ indicating a closer match to the key.
The centred form $\lam - \xstat$ standardises by the non-centrality parameter, which
matters if $\mathrm{df}$ or $\sigma$ vary across conditions.

These are not notational variants of each other. Over our runs $\sigma$ ranges from
$31.75$ to $42.59$ and $\lam$ from $856.6$ to $1541.4$, so the two scores are not
monotonically related across observations. On identical samples the forged arm
separates from the clean null with AUC $0.861$ under $-\xstat$ and $0.972$ under
$\lam - \xstat$; the genuine arm attains $1.000$ under both. A difference of that
magnitude on the forged arm, arising purely from a choice of transformation, means
the score must be treated as a pre-registered modelling decision rather than a
presentational one --- selecting between candidate scores after observing separation
on the same samples used to report it is selection on the test set. We report both
throughout and designate neither as canonical.

\subsection{Measurement harness}

The released attack script writes only the censored $p$-value to its metrics output,
so the measurement runs do not invoke it. We assemble the attack from the
repository's own primitives --- inversion, generation, and detection are all the
original implementations, unmodified --- and add a third arm the script does not
produce. Each trial yields three observations against a common watermark instance:

\begin{itemize}\itemsep2pt
\item \textbf{genuine} --- a watermarked latent generated in the target model and
      inverted back;
\item \textbf{clean} --- the same prompt generated from ordinary Gaussian noise,
      providing an empirical null against which both other arms are measured;
\item \textbf{forged} --- the genuine image inverted in the \emph{attacker} model,
      regenerated there under an unrelated prompt, and returned to the target
      detector.
\end{itemize}

The clean arm is the addition that matters. Without an empirical null there is no
basis for stating what a given score is evidence \emph{for}, only whether it crosses
a threshold.

Algorithm~\ref{alg:trial} states one trial. The two models are held on separate
devices and swapped to host memory around each other's turn, so that peak allocation
tracks the larger of the two rather than their sum.

\begin{algorithm}[t]
\small
\caption{One measurement trial. $u_T$ is the target model, $u_A$ the attacker proxy,
$\mathcal{I}$ inversion, $\mathcal{G}$ generation, $\mathcal{D}$ the Tree-Ring
detector, and $k$ the watermark key. The prompts $q_T$ and $q_A$ are unrelated.}
\label{alg:trial}
\begin{algorithmic}[1]
\State $z_T^{(w)} \gets \textsc{Watermark}(k)$ \Comment{key-derived latent}
\State $x^{(g)} \gets \mathcal{G}(u_T, q_T, z_T^{(w)})$ \Comment{genuine}
\State $x^{(c)} \gets \mathcal{G}(u_T, q_T, \varepsilon),\;
        \varepsilon \sim \mathcal{N}(0, I)$ \Comment{clean null}
\State $\hat{z}^{(w)} \gets \mathcal{I}(u_A, x^{(g)})$
        \Comment{invert in \emph{attacker} model}
\State $x^{(f)} \gets \mathcal{G}(u_A, q_A, \hat{z}^{(w)})$ \Comment{forged}
\For{$x \in \{x^{(g)}, x^{(c)}, x^{(f)}\}$}
  \State $\hat{z} \gets \mathcal{I}(u_T, x)$
  \State record $(\xstat, \lam, \mathrm{df}, \sigma, p)
         \gets \mathcal{D}(\hat{z}, k)$
\EndFor
\end{algorithmic}
\end{algorithm}

Line~4 is the entire vulnerability: the ring pattern survives inversion through a
model that never held the key, so line~5 regenerates an unrelated image from a latent
that still carries it.

We state the boundary explicitly. The attack concept, the DDIM inversion, the
Tree-Ring detector, the diffusion models and the prompt sets are the original
authors'. The loop assembly, the clean null arm, the recovery of the discarded
statistic and its validation are ours.

\section{Evaluation}
\label{sec:eval}

All results are from six trials of three arms, $n = 18$ observations, at seed $123$.

\subsection{G1 --- memory feasibility}

\begin{table}[t]
\centering\small
\caption{Peak allocated memory, measured with
\code{max\_memory\_allocated} after a single generation. Card capacity is 14.6\,GB.}
\label{tab:g1}
\begin{tabular}{lrr}
\toprule
Configuration & Peak (GB) & Load (s) \\
\midrule
SD 2.1 attacker, fp32 & 5.57  & 27.6 \\
SDXL target, fp16     & 7.36  & 73.6 \\
SDXL target, fp32     & 13.68 & 59.6 \\
\bottomrule
\end{tabular}
\end{table}

Table~\ref{tab:g1} gives the measured peaks. At full precision the pair requires
$19.25$\,GB against $14.6$\,GB available and cannot co-reside. Half precision on the
target reduces its peak by $6.32$\,GB, bringing the pair to $12.93$\,GB --- within a
single card's capacity, with headroom.

Co-residency is in fact achievable even at full precision, because the repository's
pipeline provider swaps an idle model to host memory: peak allocation tracks
$\max(\cdot)$ over the two models rather than their sum. Half precision is therefore
not strictly required, but it removes the swap from the inner loop of every trial.
\textbf{G1 passes.}

\subsection{G2 --- wall-clock cost}

A single Reprompt run completes in $338.3$\,s on a T4 against the approximately
$30$\,s the original work reports on an A40 --- a factor of $11.3$. The gap is
consistent with the difference in memory bandwidth and tensor-core generation, and it
is not prohibitive: the attack remains a sub-six-minute operation on hardware
available at no cost, which is the property that matters for a threat model.
\textbf{G2 passes.}

\subsection{G3 --- statistic recoverability and range}

All $18$ observations produced distinct values of $\xstat$, spanning $889.1$ to
$1802.2$ against $\mathrm{df} = 634$. The cross-check against the released detector
returned a maximum absolute difference of exactly $0$. The statistic is recoverable
and carries dynamic range. \textbf{G3 passes.}

The $p$-value did not behave as we had predicted. We had inferred from reading the
detector source that \code{scipy.stats.ncx2.cdf} would underflow to exactly $0.0$ on
strong genuine detections, left-censoring the genuine distribution and making
$-\log_{10} p$ undefined --- and we had used that inference to justify recovering
$\xstat$ at all. Measured: $p$ was exactly zero in $0$ of $18$ observations, with
$18$ distinct values and a smallest non-zero value of $5.45 \times 10^{-49}$. The
prediction did not reproduce. The case for recovering the raw statistic survives on
the grounds given in Section~\ref{sec:score} --- dynamic range and score
comparability --- but not on the grounds we originally gave.

\subsection{G4 --- hypothesis separation}

\begin{table}[t]
\centering\small
\caption{Per-arm outcomes ($n = 6$ each) against the values M\"uller et al.\ report
for the same model pairing. Mean $p$ is dominated by the largest observation; medians
are given for the genuine and forged arms, whose distributions span many orders of
magnitude.}
\label{tab:arms}
\begin{tabular}{lrrrr}
\toprule
Arm & $\bar{\xstat}$ & $\overline{\lam - \xstat}$ & det. & orig.\ det. \\
\midrule
genuine & $1096.5$ & $+220.0$ & $6/6$ & $1.00$ \\
clean   & $1623.9$ & $-626.3$ & $0/6$ & $0.01$ \\
forged  & $1391.2$ & $-371.3$ & $5/6$ & $0.97$ \\
\bottomrule
\end{tabular}
\end{table}

\begin{table}[t]
\centering\small
\caption{Detection $p$-values against the original work's reported means.}
\label{tab:pvals}
\begin{tabular}{lrrr}
\toprule
Arm & median $p$ & mean $p$ & orig.\ mean $p$ \\
\midrule
genuine & $8.63 \times 10^{-38}$ & $7.96 \times 10^{-24}$ & $1.10 \times 10^{-21}$ \\
clean   & $4.56 \times 10^{-1}$  & $4.80 \times 10^{-1}$  & $4.72 \times 10^{-1}$ \\
forged  & $2.30 \times 10^{-5}$  & $1.50 \times 10^{-2}$  & $5.20 \times 10^{-3}$ \\
\bottomrule
\end{tabular}
\end{table}

\paragraph{The attack reproduces.} Tables~\ref{tab:arms} and~\ref{tab:pvals} place
our measurements against the original work's reported values. The clean arm matches
closely --- mean $p$ of $0.480$ against $0.472$, and no false positives in six trials
against a reported rate of $0.01$ --- which is the arm that most directly validates
the detector as configured. The genuine arm is detected in every trial. The forged
arm crosses the $\tau = 0.0261$ threshold in five of six trials, with a median $p$
of $2.30 \times 10^{-5}$, three orders of magnitude below the threshold. A single
Reprompt run through the repository's own entry point likewise produced a genuine
$p$ of $3.53 \times 10^{-33}$ and a forged $p$ of $3.13 \times 10^{-6}$, both
detected.

\paragraph{Separation.} Against the clean null, the genuine arm separates perfectly
(AUC $1.000$ on both scores). The forged arm separates at AUC $0.861$ on $-\xstat$
and $0.972$ on $\lam - \xstat$ (Section~\ref{sec:score}). \textbf{G4 passes.}

\begin{figure}[t]
\centering
\includegraphics[width=\columnwidth]{figures/gate4_separation.png}
\caption{Left: distribution of the evidence score $-\xstat$ by arm. Right: per-trial
values of $\xstat$. The forged arm occupies an intermediate band that overlaps the
clean null at its upper edge.}
\label{fig:sep}
\end{figure}

\paragraph{A divergence worth naming.} Our forged detection rate of $5/6 = 0.83$ sits
below the $0.97$ reported originally. At $n = 6$ this is not a statistically
meaningful difference --- a $95\%$ Wilson interval on $5/6$ spans roughly $0.44$ to
$0.97$ --- and we do not claim the attack is weaker than reported. But the shape of
the failure is informative. Figure~\ref{fig:sep} shows the forged arm occupying a
band between the clean null and the genuine arm, and the single undetected forgery
has $\xstat = 1558.3$, inside the clean range of $1483.7$ to $1802.2$. The forgery
did not fail marginally; it produced a statistic indistinguishable from an
unwatermarked image.

This is the substantive observation of the section. A detector operating as a binary
test reports the five successful forgeries and the six genuine images identically,
while the intermediate position of the forged distribution --- clearly visible in the
raw statistic and entirely absent from the thresholded decision --- is exactly the
information that would distinguish them. The detector computes that information and
discards it.

\section{Discussion}
\label{sec:disc}

\subsection{Where reproduction effort goes}

Invoking the released attack is a single line of code. Our reproduction notebook is
$375$ lines. Table~\ref{tab:effort} shows the distribution, and the ratio is worth
stating plainly: the published artifact is not the expensive part of reproducing a
published artifact.

\begin{table}[t]
\centering\small
\caption{Distribution of reproduction code by purpose.}
\label{tab:effort}
\begin{tabular}{lrl}
\toprule
Component & Lines & Category \\
\midrule
Environment repair, provenance & 132 & portability \\
Memory measurement (G1)        & 36  & instrumentation \\
Invoking the released attack   & 1   & original artifact \\
Timing wrapper (G2)            & 22  & instrumentation \\
Precision localisation         & 49  & portability \\
Statistic recovery and harness & 71  & extension \\
Analysis and gate verdicts     & 64  & extension \\
\bottomrule
\end{tabular}
\end{table}

This is not a criticism of the artifact, which ran and produced results consistent
with its paper. It is an observation about what artifact evaluation measures.
Executability is verified in the environment the authors specify; portability across
hardware generations, driver stacks and library versions is a distinct property, and
it is the one that determines whether the work can be built upon a year later on
different hardware. Roughly $45\%$ of our code exists solely to establish an
environment in which the remaining $55\%$ can run.

\subsection{Detector interfaces determine what can be asked}

The finding we consider most transferable is that the detector's return type
constrains the questions downstream work can pose. Returning a thresholded decision
or a CDF value is sufficient for deployment and insufficient for analysis. Our entire
instrumentation effort exists to undo one line of the released detector --- the
decision to return $p$ rather than the tuple $(\xstat, \lam, \mathrm{df}, \sigma)$
from which $p$ is computed.

The cost of exposing those quantities is negligible; they are already computed. The
cost of not exposing them, for a downstream consumer, is a reimplementation that must
then be validated against the original. We suggest that detectors intended for
research use return their sufficient statistics alongside any derived decision. This
is an argument about interface design, not about the correctness of the scheme or the
admissibility of any particular evidence.

\subsection{Two predictions that did not survive measurement}

We reported two inferences drawn from source inspection that our own measurements
contradicted, and we state them together because they share a mechanism.

The first was that the detector's $p$-value is left-censored at zero. The reasoning
was not unreasonable: the CDF of a non-central $\chi^2$ evaluated far into its lower
tail is a plausible underflow candidate, and the released code prints $p$-values as
\code{0.0} in its console output, which we took as corroboration. Both the reasoning
and the apparent corroboration were wrong. The printed zeros were a display-precision
artifact; the underlying values were finite and distinct across all eighteen
observations, down to $5.45 \times 10^{-49}$.

The second was that half-precision detection failure localises to the autoencoder
(Section~\ref{sec:fp16}). The probe we built to confirm it returned bit-identical
$p$-values under both dtypes and no non-finite values in either stage. The failure we
had observed came from a different execution path, and the attribution was
unsupported.

A third, smaller correction: we initially attributed a PSNR calculation fix to commit
\code{43eec7e}, which modifies documentation only. The fix is \code{d5885a3}.

What the three share is that each was inferred from reading rather than from running,
and that in two of the three cases an artifact of \emph{presentation} --- console
formatting, sane-looking quality metrics --- was mistaken for evidence about
underlying values. We report them because the same inferential shortcut is available
at every point in a reproduction, and because a study that reports only the
predictions that held has not demonstrated that it was capable of detecting one that
did not.

\subsection{Implications for the attack}

Our reproduction supports the original conclusion. Forgery succeeded in five of six
trials on the first configuration we tried, at $338.3$\,s per attempt on free
hardware, using a proxy model that shares no weights with the target. The attacker
requirement is a single watermarked reference image, and the median forged $p$-value
sits three orders of magnitude below the calibrated detection threshold.

The property this exploits is specific rather than general. Tree-Ring detection is a
similarity test against a recoverable artifact, and the secret used to verify is
identical to the secret used to inject. Any scheme with both properties admits the
same attack shape, and the transferability of the pattern across models means the
attacker need not possess the target model at all. This is not a general law relating
robustness to forgeability --- schemes that verify with a secret distinct from the one
they inject with are not implicated --- but it does mean that robustness evaluations
of this family of schemes measure a property that is orthogonal to the one a
provenance claim depends on.

\section{Limitations}
\label{sec:limits}

We state these plainly; several materially constrain what our numbers support.

\textbf{A single watermark key.} All observations share one Tree-Ring seed
($\code{w\_seed} = 999999$). The variation we report is therefore within-key. These
are not independent draws over the key space, and the separation in
Section~\ref{sec:eval} must not be read as a distribution over watermarks. This is
the most serious limitation of the present study, and it is the first thing corrected
in subsequent work.

\textbf{Sample size.} Eighteen observations from six trials, six per arm. This is
sufficient to establish that the attack reproduces and that the statistic is
recoverable; it is insufficient for any distributional claim. The AUC values in
Section~\ref{sec:eval} should be read as ordering indicators rather than calibrated
estimates, and the gap between our forged detection rate of $5/6$ and the reported
$0.97$ is not resolvable at this sample size.

\textbf{Attack scope.} We reproduce Reprompt only. The optimisation-based Imprint
variant costs roughly two orders of magnitude more per run in the original authors'
own reporting, which places a replication with meaningful sample size outside our
compute allocation. We make no claim about Imprint's behaviour on our hardware.

\textbf{Precision divergence.} The original results are fp32 throughout; ours are
mixed precision. Our probe found fp16 and fp32 detection to be bit-identical on the
target (Section~\ref{sec:fp16}), but on a single image and most likely because the
library upcast the autoencoder irrespective of our request. We have not established
that fp16 is numerically neutral for the detection statistic in general.

\textbf{Excluded target models.} FLUX.1-dev requires \code{bfloat16}, unavailable on
the T4. PixArt-$\Sigma$ is supported by the released code but runs a DPM scheduler
rather than DDIM, and the original work's own ablation shows the scheduler
substantially affects latent recovery quality --- so an SDXL/PixArt comparison would
confound architecture with scheduler. We report neither.

\textbf{Pristine images only.} All observations are of unprocessed images. Realistic
handling --- recompression, resizing, cropping, screenshotting --- is untested here,
and both genuine and forged populations would need to pass through matched channels
before any conclusion about behaviour under such handling could be drawn.

\textbf{Attacker model substitution.} The original SD~2.1 repository was withdrawn by
its publisher during this work and we used a mirror. We verified the weights load and
produce the expected attack behaviour, but did not perform a bitwise comparison
against the original release.

\section{Conclusion}

We reproduced the Reprompt semantic watermark forgery attack against Tree-Ring on
Stable Diffusion~XL, using the authors' released implementation on free-tier hardware
providing roughly a third of the memory of the original evaluation environment. The
attack reproduces: genuine images are detected in six of six trials, clean images in
none, and forged images in five of six, at $338.3$\,s per attempt against the $30$\,s
reported on datacentre hardware. The clean arm matches the original work's reported
behaviour closely (mean $p$ of $0.480$ against $0.472$), which is the arm that most
directly validates the detector as configured.

The result we consider most transferable concerns the detector's interface rather
than the attack. The released detector computes a non-central $\chi^2$ statistic and
returns only its CDF. Recovering that statistic --- bitwise identical to the
detector's own computation --- shows that two natural scores derived from it separate
the forged arm at AUC $0.861$ and $0.972$ on identical samples, so the choice of
score is a modelling decision that must be fixed before observing separation. It also
shows the forged arm occupying an intermediate band between the clean null and the
genuine arm: information the detector computes, then discards at the point of
returning a decision.

We also report predictions that failed. We had inferred that the detector's
$p$-value was left-censored at zero; it was never zero across eighteen observations,
with a minimum of $5.45 \times 10^{-49}$. We had inferred that half-precision
detection failure localised to the autoencoder; the probe built to confirm it
returned bit-identical results under both dtypes. In each case an artifact of
presentation was mistaken for evidence about underlying values. We report them
because a reproduction that surfaces only its confirmed predictions has not
demonstrated that it could detect a failed one.

Immediate future work varies the watermark key across clusters, which is the binding
limitation of the present study, and passes both genuine and forged populations
through matched laundering channels so that the distributions being compared have
undergone identical handling.

\section*{Division of Work}

Saifur Rahman Tamim designed the gate protocol, implemented the environment
reconstruction and provenance pinning, designed and ran the precision probe, and
implemented the statistic recovery and its validation. Md Taslimul Hasan Toufique
\NUM{contribution}. Fahim Ibney Hafiz \NUM{contribution}. A.M. Tayeful Islam advised
on experimental design and reviewed the manuscript. All authors contributed to the
analysis and to the writing.

\section*{Availability}

The reproduction notebook, gate artifacts and pinned fork commit are available at
\NUM{repository URL or ``on request''}.

\balance
\bibliographystyle{plainnat}
\bibliography{refs}
\end{document}
